%!TeX root=MemoriaTFG.tex

\chapter{Discusión}
\label{cap:discusion}

Este capítulo lee en conjunto los resultados del
capítulo~\ref{cap:resultados} y del capítulo sobre Dermapixel. Primero
los sitúa frente a la literatura previa (qué confirman y qué matizan);
después destila los hilos transversales que de verdad importan; y, por
último, declara con honestidad hasta dónde llegan ---y hasta dónde no---
las conclusiones.

\section{Comparación con trabajos previos}

\paragraph{Reproducción de PanDerm.} Las cifras de linear probe sobre
los datasets compartidos (tabla~\ref{tab:lp-panderm}) caen dentro del
rango que reportan los autores de PanDerm~\cite{panderm2025} y conservan
la misma relación entre PanDerm Base y Large, pese a ejecutarse en una
plataforma distinta de la original (arquitectura aarch64, BF16 sobre
NVIDIA Grace Hopper GB10). Que la ventaja de PanDerm Large se mantenga
sobre datasets externos, ajenos a su desarrollo, aporta evidencia
independiente de que generaliza: no depende de los conjuntos con los que
se entrenó.

\paragraph{Segmentación.} En segmentación binaria de lesión, el modelo
adaptado iguala o supera ligeramente lo publicado para PanDerm. Pero el
mensaje de fondo matiza la lectura habitual: lo que decide el resultado
no es qué segmentador se use ---las arquitecturas se separan apenas
$\sim 1$\,pp de Dice--- sino dar al modelo la localización de la lesión
(la caja). De hecho, el modelo coincide con el consenso de varios
dermatólogos ($0{,}915$ de Dice) más de lo que esos dermatólogos
coinciden entre sí ($0{,}728$): la tarea está, sobre este dataset, cerca
de su techo práctico.

\paragraph{Zero-shot multimodal.} DermLIP alcanza el AUROC que reporta
la literatura sobre Derm1M~\cite{derm1m} solo cuando el texto se formula
como en su desarrollo. La contribución aquí es cuantificar lo frágil que
es eso: sobre el mismo encoder visual, cambiar el tokenizer y la
redacción de los prompts mueve el AUROC de forma notable. La mayoría de
trabajos reportan una única configuración textual; medir esa
sensibilidad de forma sistemática invita a una cautela metodológica:
diferencias atribuidas a la arquitectura pueden venir, en realidad, de
cómo se escribió la evaluación.

\paragraph{Disparidad por fototipo.} El descenso de rendimiento sobre
los fototipos más oscuros confirma lo descrito en
Fitzpatrick17k~\cite{fitzpatrick17k} y la motivación de DDI~\cite{ddi}:
el sesgo por representación desigual de tonos de piel sigue abierto, aun
en los modelos fundacionales dermatológicos más recientes. Lo nuevo es
compararlo de forma homogénea entre cinco encoders bajo un mismo
protocolo, lo que revela que rendimiento global y equidad no van de la
mano, como se detalla en la sección siguiente.

\section{Resultados alcanzados}

\paragraph{No hay un método óptimo único: especialización por tarea.}
\label{sec:disc-especializacion}
El hilo que atraviesa todo el trabajo es que \textbf{ningún modelo ni
ninguna receta gana en todo}. La mejor opción depende a la vez de la
tarea (clasificar, segmentar, recuperar casos), del nivel de detalle
diagnóstico (L1/L2/L3) y de la métrica que se mire. De ahí salen dos
lecciones prácticas. La primera: con pocos datos etiquetados conviene
una adaptación ligera ---LoRA (PEFT) rinde mejor que el fine-tuning
denso entrenando muchísimos menos parámetros---, mientras que con un
dataset grande el fine-tuning sí compensa.\label{sec:disc-lp-vs-ft} La
segunda: clasificar y recuperar casos no son el mismo objetivo, y
optimizar para uno puede empeorar el otro. La consecuencia de
diseño es preferir un sistema modular, con el componente adecuado para
cada tarea, en lugar de buscar un único modelo que lo haga todo.

\paragraph{Cómo de fiables son estas comparaciones.}
\label{sec:disc-cv-robustez}
Sobre un dataset pequeño, una sola partición engaña. La validación
cruzada por caso sobre DermapixelAI~1.0 (tabla~\ref{tab:dermapixel-cv})
lo enseña: la cifra de una partición fija puede ser bastante optimista en
métricas con umbral (como la BAcc), mientras que la AUROC apenas cambia.
Por eso las jerarquías entre métodos en L2 y L3 se leen como indicios
sujetos a incertidumbre, no como rankings cerrados, y se anclan en la
estimación cruzada, no en el número de una partición afortunada.

\paragraph{La familia PanDerm bate a los generalistas, pero con cierre
estadístico parcial.}
\label{sec:disc-llm}
Los modelos fundacionales dermatológicos superan con claridad a los
grandes modelos de lenguaje multimodales generalistas: sobre HAM10000 la
distancia entre el mejor especialista adaptado y el mejor LLM multimodal
en zero-shot es de varias decenas de puntos de accuracy. En dermatología
visual, el conocimiento del dominio pesa más que la escala bruta del
modelo; además, el especialista es más barato (se ejecuta en local). Esa
superioridad «dominio $>$ generalista» se contrastó formalmente con el
test de DeLong sobre tres endpoints de malignidad
(sección~\ref{sec:equidad-delong}), pero solo Fitzpatrick17k tiene
potencia muestral suficiente; en HAM10000 y DermapixelAI~1.0 ningún
contraste alcanza significación. Importa subrayarlo con honestidad:
\textbf{falta de potencia no es lo mismo que ausencia de efecto}. El
cierre estadístico respalda la jerarquía frente a los generalistas, no
una ordenación fina entre los dos líderes específicos (PanDerm y
DermLIP).

\paragraph{Hacen falta pocas etiquetas ---a veces ninguna.}
\label{sec:disc-zs}
Una consecuencia práctica de que el encoder ya esté maduro es que casi no
hacen falta etiquetas para aprovecharlo, y el etiquetado experto es el
verdadero cuello de botella en dermatología. La AUROC satura muy pronto:
con el $1\%$ de HAM10000 ($82$ imágenes) el linear probe ya alcanza
$0{,}918$, y en PAD-UFES-20 bastan $14$ imágenes para recuperar el
$95{,}5\%$ de la AUROC final ---lo que sigue mejorando con el volumen es
la balanced accuracy, más sensible a las clases minoritarias---. En el
extremo, ni siquiera hacen falta etiquetas: en las subcategorías L2 de
DermapixelAI~1.0, DermLIP v2 en zero-shot iguala (AUROC $0{,}794$) al
mejor linear probe supervisado ($0{,}793$). El reverso es que ese régimen
zero-shot es muy sensible a cómo se redacta el texto ---el tokenizer y los
prompts mueven el AUROC de forma notable, y pasar los prompts a castellano
penaliza---, lo que sitúa el factor limitante en el lenguaje, no en la
imagen, y justifica una rama textual en español.

\paragraph{El cuello de botella ya no es el modelo, es el dato.} En
varios experimentos, tocar el diseño del sistema ---el texto de los
prompts, la estrategia de adaptación, la ontología, el equilibrio de
clases, la recuperación por similitud--- mueve el rendimiento tanto o
más que cambiar el encoder visual. Esto es coherente con que los
encoders ya han madurado: alcanzada cierta calidad de las
representaciones, el margen de mejora se desplaza de la arquitectura
hacia los datos (su calidad, diversidad y cobertura) y la long tail. En
el nivel diagnóstico fino (L3), de hecho, la limitación deja de ser el
modelo y pasa a ser la escasez de ejemplos por clase.

\paragraph{Equidad por fototipo.}
\label{sec:disc-equidad}
Los cinco encoders rinden peor sobre los fototipos oscuros, y la lectura
honesta exige mirar la equidad junto al rendimiento global: la paradoja
de BiomedCLIP ---el menor gap entre fototipos pero la peor BAcc
global--- muestra que \textbf{un modelo malo por igual en todos los
subgrupos parece equitativo sin serlo}. Bajo esa lectura conjunta,
PanDerm Large ofrece el mejor compromiso, y la validación
estratificada por fototipo es un requisito de seguridad, no un extra.

\paragraph{Por dentro, el modelo ya usa los conceptos del dermatólogo.}
Más allá de acertar, para que un especialista confíe en el sistema hace
falta saber \emph{por qué} decide. El Sparse Autoencoder muestra que las
representaciones de PanDerm ya codifican, sin que nadie se lo enseñe, los
conceptos clásicos de la dermatoscopia: las mejores \emph{features}
detectan criterios del \emph{Seven-Point Checklist} con AUROC de $0{,}68$
a $0{,}82$ ---retículo pigmentado, velo azul-blanquecino, estructuras
vasculares, estrías--- y cada concepto se reparte en un grupo coherente de
\emph{features}, no en una única neurona. Forzar la decisión a pasar solo
por esos siete conceptos (un \emph{Concept Bottleneck Model}) cuesta algo
de precisión (AUROC $0{,}890 \to 0{,}832$) a cambio de una decisión
enteramente trazable, y el modelo pondera más alto justo los criterios que
el propio checklist considera más diagnósticos. Lo que de verdad abre esta
línea no es replicar escalas en inglés, sino construir un vocabulario
propio en castellano con la Dra.~Taberner: de $16$ conceptos propuestos, el
SAE ya captura $7$ ---incluido uno con solo $15$ casos positivos---, y los
que no figuran en las escalas clásicas son precisamente los que justifican
seguir.

\section{Limitaciones}
\label{cap:limitaciones}

Las conclusiones anteriores se sostienen como patrones cualitativos
robustos, pero conviene acotar su dominio de validez. Se agrupan aquí,
en frases cortas, las limitaciones del trabajo.

\paragraph{Disponibilidad.}
\label{sec:lim-pesos}
Los pesos de DermFM-Zero no son públicos a la fecha de cierre, de modo
que solo se trata conceptualmente y no se evalúa. Los modelos cerrados
accesibles por API (GPT-4o, Gemini) no permiten fijar la versión exacta
ni el posproceso del proveedor, así que sus cifras son condicionales a la
fecha de evaluación. Y los conjuntos de preentrenamiento internos
(PanDerm, Derm1M) no se liberan: la composición declarada de esos datos
descansa en lo que dicen sus autores, no en una auditoría independiente.
Buena parte del trabajo se apoya, además, en modelos de terceros, con los
sesgos de diseño que eso herede.

\paragraph{Datos.}
\label{sec:lim-leakage}
\label{sec:lim-muestrales}
Los datasets externos proceden sobre todo de EE.\,UU., Australia, Europa
central y Brasil, y las ramas de texto se preentrenan en inglés; la
población hispanohablante apenas está representada (brecha que
DermapixelAI~1.0 empieza a cubrir). Las taxonomías nativas de cada
dataset son incompatibles entre sí: la ontología L1/L2/L3 lo mitiga pero
no lo elimina, y las comparaciones a nivel L3 siguen siendo delicadas. La
auditoría de solapamiento con Derm1M halla un único dataset contaminado
(Dermnet), y solo detecta duplicados exactos por hash, no recortes ni
reescalados. Por último, los tamaños por subgrupo del experimento de
equidad son asimétricos ---el fototipo~VI apenas reúne $73$ casos de
test---, lo que ensancha los intervalos y limita la precisión del gap,
aunque no su sentido cualitativo.

\paragraph{Diseño.}
\label{sec:lim-semilla}
\label{sec:lim-estadistica}
Varios protocolos (linear probe, eficiencia en etiquetas, equidad)
corren con una sola semilla, así que no reportan su variabilidad; se
asume porque la magnitud de los efectos clave excede con holgura esa
variabilidad esperable, y la segmentación sí incluye validación cruzada.
El cierre estadístico con DeLong y la corrección de Holm tiene potencia
únicamente en Fitzpatrick17k; en el resto, las diferencias se leen como
descriptivas. Y tanto el zero-shot como los LLMs son sensibles a cómo se
redacta el prompt, por lo que esas cifras son condicionales a las
formulaciones evaluadas y deben leerse como cota, no como techo.

\paragraph{Dataset DermapixelAI 1.0.} El dataset está muy desbalanceado
por modalidad (predomina con mucho la fotografía clínica sobre la
dermatoscópica), lo que lo hace poco apto para benchmarks
dermatoscópicos. Su cobertura de la ontología L3 es parcial y arrastra
una long tail marcada. Sus tres cifras de validación ---mapeo ontológico,
revisión visual de la Dra.~Taberner y coherencia del texto--- miden cosas
distintas y no son intercambiables. Y la partición case-aware, correcta
para evitar fuga entre imágenes del mismo caso, deja un test reducido;
por eso las comparaciones sobre este dataset se anclan en la validación
cruzada por caso y no en la cifra de una partición fija.

\paragraph{Generalización clínica.}
\label{sec:alcance}
El trabajo mide rendimiento sobre fotografías de archivos curados o
clínicos, no sobre el flujo real de teledermatología (dispositivos
heterogéneos, iluminación y encuadre sin estandarizar): extrapolar a ese
escenario exige validación específica. No hay validación prospectiva
sobre cohorte hospitalaria con aprobación de un comité ético, por
estructura del propio TFG. Y quedan fuera, por decisión de alcance,
varios elementos que conviene reconocer: el preentrenamiento desde cero
de un modelo propio, la integración con el PACS del hospital, el
reentrenamiento multilingüe a gran escala, la certificación CE como
producto sanitario y el análisis sistemático de calibración y de utilidad
clínica. Nada de esto invalida los patrones cualitativos del trabajo;
acota dónde son válidos y marca la continuación natural hacia un
despliegue real.
